% Part: turing-machines
% Chapter: undecidability
% Section: representing-tms

\documentclass[../../../include/open-logic-section]{subfiles}

\begin{document}

\olfileid{tur}{und}{rep}
\olsection{ٹیورنگ مشیناں دی منطقی نمائندگی}

\begin{explain}
ٹیورنگ مشیناں تے اوہناں دے چلن نوں پہلے درجے دی منطق دے
اک !!a{sentence} نال ظاہر کرن لئی سانوں اک موزوں زبان دی
تعریف کرنی پئے گی۔ ایس زبان دے دو حصے نیں: مشین دیاں مکمل
حالتاں بیان کرن لئی !!{predicate}s، تے چال دے قدماں (“ویلیاں”)
تے پٹی اُتے تھاماں نوں نمبر دین لئی عبارتاں۔

اسی دو قسماں دے !!{predicate}s ورتاں گے؛ دونیں دو-جگہی
نیں: ہر حالت~$q$ لئی اک !!a{predicate}~$\Obj Q_q$، تے پٹی دے
ہر نشان~$\sigma$ لئی اک !!a{predicate}~$\Obj S_\sigma$۔ پہلے
والے سانوں~$M$ دی حالت تے اوہدے پڑھن لکھن والے سر دی تھاں
دسن دیندے نیں؛ دوجے والے دسدے نیں کہ پٹی اُتے کیہڑے نشان نیں۔

پٹی اُتے سر دیاں تھاماں تے پورے ہو چکّے قدماں دی گنتی ظاہر
کرن لئی سانوں عدد ظاہر کرن دا طریقہ چاہیدا اے۔ ایس لئی
اسی اک !!a{constant}~$\Obj 0$ تے اک $1$-جگہی فنکشن~$\prime$،
یعنی جانشین فنکشن، ورتدے آں۔ رواج مطابق ایہ اپنی دلیل
توں \emph{بعد} لکھیا جاندا اے (تے اسی بریکٹ چھڈ دیندے آں)۔

ہر عدد~$n$ لئی اک مقررہ ٹرم~$\num{n}$ اے، یعنی~$n$ دی
\emph{عددی عبارت}، جیڑی~$\Lang L_M$ وچ اوہنوں ظاہر کردی اے۔
$\num{0}$، $\Obj 0$ اے؛ $\num{1}$، $\Obj 0'$ اے؛
$\num{2}$، $\Obj 0''$ اے؛ تے ایداں اگے۔ زیادہ رسمی طور اُتے:
\begin{align*}
\num{0} & = \Obj 0 \\
\num{n+1} &= \num{n}'
\end{align*}
ٹرم~$\num{0}$، یعنی~$\Obj 0$، پٹی دے سب توں کھبے خانے
دا ناں وی اے تے پہلا قدم چُکن توں پہلاں دے ویلے
(ابتدائی مکمل حالت) دا ناں وی۔ ٹرم~$\num{1}$، یعنی~$\Obj 0'$،
سب توں کھبے خانے دے سجّے پاسے والے خانے دا ناں اے، تے
پہلا قدم چُکن توں بعد دے ویلے دا وی؛ ایداں ہی اگے۔

اسی پٹی اُتے تھاماں دی ترتیب (جدوں ایس دا مطلب “دے کھبے
پاسے” ہووے) تے چلن دے قدماں دی ترتیب (جدوں مطلب “توں پہلاں”
ہووے) دونیں ظاہر کرن لئی اک !!a{predicate}~$<$ وی ورتاں گے۔

زبان مقرر کرن توں بعد اسی~$!T(M,
w)$ دے “مسلّمے” لکھاں گے، یعنی اوہ !!{sentence}s جیہڑے
اکٹھے ہو کے مشین~$M$ دا انپٹ~$w$ اُتے چلن بیان
کردے نیں۔ کجھ !!{sentence}s، $\Obj 0$، $\prime$ تے~$<$
دیاں شرطاں مقرر کرن گے؛ کجھ !!{sentence}s ابتدائی مکمل
حالت بیان کرن گے؛ تے کجھ !!{sentence}s بیان کرن گے کہ
مشین~$M$ کوئی خاص حکم چلاؤن توں بعد کس مکمل حالت وچ اے۔
\end{explain}

\begin{defn}
  \ollabel{defn:tm-descr}
جے ٹیورنگ مشین $M = \tuple{Q, \Sigma, q_0, \delta}$ دتی
ہووے، تاں زبان~$\Lang L_M$ وچ ایہ چیزاں ہوندیاں نیں:
\begin{enumerate}
\item اک دو-جگہی !!{predicate} $\Obj Q_q(x, y)$ ہر حالت~$q \in
  Q$ لئی۔ مطلب ایہ
  اے کہ $\Obj Q_q(\num{m}, \num{n})$ آکھدا اے: “$n$
  قدم چُکن توں بعد~$M$ حالت~$q$ وچ اے تے~$m$ویں خانہ
  پڑھ رہی اے۔”
\item اک دو-جگہی !!{predicate} $\Obj S_\sigma(x, y)$ ہر
  نشان~$\sigma\in \Sigma$ لئی۔ مطلب ایہ اے کہ $\Obj S_\sigma(\num{m},
  \num{n})$ آکھدا اے: “$n$ قدم چُکن توں بعد~$m$ویں خانے
  وچ نشان~$\sigma$ اے۔”
\item اک !!{constant} $\Obj 0$
\item اک یک-مقامی !!{function} $\prime$
\item اک دو-جگہی !!{predicate} $<$
\end{enumerate}
\end{defn}

ٹیورنگ مشین~$M$ دا انپٹ $w = \sigma_{i_1}\dots\sigma_{i_k}$
اُتے چلن بیان کرن والے !!{sentence}s ایہ نیں:
\begin{enumerate}
\item عدداں تے~$<$ بارے مسلّمے:
\begin{enumerate}
\item اک !!^a{sentence} جیڑا آکھدا اے کہ ہر عدد اپنے جانشین توں
  گھٹ اے:
\[
\lforall[x][x < x']
\]
\item اک !!^a{sentence} جیڑا یقینی بناندا اے کہ~$<$ متعدی اے:
\[
\lforall[x][\lforall[y][\lforall[z][
      ((x < y \land y < z) \lif x < z)]]]
\]
\end{enumerate}
\item انپٹ دی ابتدائی مکمل حالت بیان کرن والے مسلّمے:
\begin{enumerate}
\item $0$ قدم چُکن توں بعد، یعنی مشین چلن توں پہلاں،
  $M$ ابتدائی حالت~$q_0$ وچ اے تے خانہ~$1$ پڑھ رہی اے:
\[
\Obj Q_{q_0}(\num{1}, \num{0})
\]
\item پہلے $k+1$ خانیاں وچ ترتیب وار نشان $\TMendtape$،
  $\sigma_{i_1}$، \dots، $\sigma_{i_k}$ نیں:
\[
\Obj S_\TMendtape(\num{0}, \num{0}) \land
\Obj S_{\sigma_{i_1}}(\num{1}, \num{0}) \land
\dots \land
\Obj S_{\sigma_{i_k}}(\num{k}, \num{0})
\]
\item باقی پٹی خالی اے:
\[
\lforall[x][(\num{k} < x \lif \Obj S_\TMblank(x, \num{0}))]
\]
\end{enumerate}
\item اک مکمل حالت توں اگلی مکمل حالت وچ تبدیلی بیان کرن
  والے مسلّمے:

اگلے حصے لئی~$!A(x, y)$ نوں اوہ عطف منّو جیہڑا ایس روپ
والے سارے !!{sentence}s جوڑ کے بندا اے:
\[
\lforall[z][
  (((z < x \lor x < z) \land \Obj S_\sigma(z, y))
  \lif \Obj S_\sigma(z, y'))]
\]
جتھے~$\sigma \in \Sigma$ اے۔ اسی~$!A(\num{m},\num{n})$
ایہ آکھن لئی ورتدے آں: “خانہ~$m$ توں علاوہ پٹی
$n+1$ قدم چُکن توں بعد اوہو جیہی اے جیہی~$n$ قدم
چُکن توں بعد سی۔”
\begin{enumerate}
\item \ollabel{rep-right} ہر حکم $\delta(q_i, \sigma) =
  \tuple{q_j, \sigma', \TMright}$ لئی ایہ !!{sentence}:
\begin{align*}
& \lforall[x][\lforall[y][(
   (\Obj Q_{q_i}(x, y) \land \Obj S_{\sigma}(x, y)) \lif {}]] \\
&\qquad   (\Obj Q_{q_j}(x', y') \land \Obj S_{\sigma'}(x, y') \land
!A(x, y)))
\end{align*}
ایہ آکھدا اے: جے~$y$ قدم چُکن توں بعد مشین حالت~$q_i$
وچ خانے~$x$ نوں پڑھ رہی اے، جتھے نشان~$\sigma$ اے، تاں
$y+1$ قدم چُکن توں بعد اوہ خانہ~$x+1$ پڑھ رہی اے،
حالت~$q_j$ وچ اے، خانے~$x$ وچ ہن~$\sigma'$ اے، تے
خانہ~$x$ توں علاوہ ہر خانے وچ اوہو نشان اے جیڑا
$y$ قدم چُکن توں بعد سی۔

\item \ollabel{rep-left} ہر حکم $\delta(q_i, \sigma) =
  \tuple{q_j, \sigma', \TMleft}$ لئی ایہ !!{sentence}:
% PNB-CMP-030: the old scanned/written square is x', so the frame condition exempts x', not x.
\begin{align*}
& \lforall[x][\lforall[y][
    ((\Obj Q_{q_i}(x', y) \land \Obj S_{\sigma}(x', y)) \lif {}]]\\
& \qquad   (\Obj Q_{q_j}(x, y') \land \Obj S_{\sigma'}(x', y') \land
!A(x', y))) \land {}\\
& \lforall[y][((\Obj Q_{q_i}(\num{0}, y) \land \Obj S_{\sigma}(\num{0},
    y)) \lif {}]\\
& \qquad (\Obj Q_{q_j}(\num{0}, y') \land \Obj S_{\sigma'}(\num{0},
  y') \land !A(\num{0}, y)))
\end{align*}
ذرا سوچو ایہ کم کیویں کردا اے: ہن اسی “جے مشین خانہ~$x$
پڑھ رہی ہووے \dots” توں نہیں، بلکہ “جے مشین خانہ~$x+1$
پڑھ رہی ہووے \dots” توں شروع کردے آں۔ کھبے ول جان دا
مطلب اے کہ اگلے قدم وچ مشین خانہ~$x$ پڑھدی اے؛ پر
نشان خانہ~$x+1$ وچ لکھدی اے۔ اسی ایویں ایس لئی کردے
آں کہ ساڈے کول تفریق یا پیشرو فنکشن دا نشان نہیں۔

دھیان دیو کہ $x+1$ روپ دے عدد $1$، $2$، \dots{} نیں؛
ایس وچ اوہ صورت شامل نہیں جدوں مشین خانہ~$0$ پڑھ رہی
ہووے تے اوہنوں کھبے ول جانا ہووے (ایہ کر نہیں سکدی،
ایس لئی اپنی تھاں رہندی اے)۔ اوہ خاص صورت عطف دے دوجے
جزو نال بیان ہوئی اے: جے~$y$ قدم چُکن توں بعد مشین
خانہ~$0$ حالت~$q_i$ وچ پڑھ رہی اے تے خانہ~$0$
وچ نشان~$\sigma$ اے، تاں~$y+1$ قدم چُکن توں بعد
اوہ ہن وی خانہ~$0$ پڑھ رہی اے، ہن حالت~$q_j$
وچ اے، خانے~$0$ دا نشان~$\sigma'$ اے، تے
خانہ~$0$ توں علاوہ خانے اوہو نشان رکھدے نیں
جیہڑے~$y$ قدم چُکن توں بعد اوہناں وچ سن۔
\item \ollabel{rep-stay} ہر حکم $\delta(q_i, \sigma) =
  \tuple{q_j, \sigma', \TMstay}$ لئی ایہ !!{sentence}:
\begin{align*}
& \lforall[x][\lforall[y][(
   (\Obj Q_{q_i}(x, y) \land \Obj S_{\sigma}(x, y)) \lif {}]] \\
&\qquad   (\Obj Q_{q_j}(x, y') \land \Obj S_{\sigma'}(x, y') \land
!A(x, y)))
\end{align*}
\end{enumerate}
\end{enumerate}
اُتلے سارے !!{sentence}s جوڑ کے بنے عطف نوں~$!T(M, w)$
آکھو، جیہڑا ٹیورنگ مشین~$M$ تے انپٹ~$w$ لئی اے۔

ایہ ظاہر کرن لئی کہ~$M$ آخرکار رُکدی اے، سانوں
اک ایہو جیہا !!{sentence} لبھنا پئے گا جیڑا آکھے: “کجھ
قدم چُکن توں بعد حالتی تبدیلی دا فنکشن غیر تعریف
شدہ ہووے گا۔” $X$~اوہناں ساریاں جوڑیاں
$\tuple{q, \sigma}$ دا سیٹ ہووے جِنہاں لئی
~$\delta(q, \sigma)$ غیر تعریف شدہ اے۔ فیر
$!E(M, w)$ ایہ !!{sentence} ہووے:
\[
\lexists[x][\lexists[y][(\bigvee_{\tuple{q, \sigma} \in
      X}(\Obj Q_q(x, y) \land \Obj S_\sigma(x, y)))]]
\]

جے اسی مقررہ رُکّن والی حالت~$h$ رکھنے والی ٹیورنگ
مشین ورتیے، تاں ایہ ہور وی سوکھا اے: فیر
!!{sentence}~$!E(M, w)$
\[
\lexists[x][\lexists[y][\Obj Q_h(x, y)]]
\]
ایہ ظاہر کردا اے کہ مشین آخرکار رُکدی اے۔

\begin{prop}
\ollabel{prop:mlessk}
جے $m < k$ ہووے، تاں $!T(M, w) \Entails \num{m} < \num{k}$۔
\end{prop}

\begin{proof}
مشق۔
\end{proof}

\begin{prob}
\olref[tur][und][rep]{prop:mlessk} ثابت کرو۔
(اشارہ: $k-m$ اُتے استقرا ورتو۔)
\end{prob}

\end{document}
